\chapter{Oliguria}

\section*{Definition}
Urine output less than 400~mL per day in adults (less than 0.5~mL/kg/hour in children), indicating decreased renal perfusion, intrinsic renal injury, or post-renal obstruction.

\section*{Pathophysiology}
Oliguria arises from one of three mechanisms: prerenal hypoperfusion activates the renin-angiotensin-aldosterone system and antidiuretic hormone to conserve water and sodium; intrinsic renal damage (acute tubular necrosis, glomerulonephritis) disrupts nephron function; or post-renal obstruction increases back-pressure that reduces glomerular filtration pressure.

\section*{Etiology}
\begin{itemize}
  \item \textbf{Pre-renal:} Hypovolemia (hemorrhage, dehydration, vomiting, diarrhea, burns, diuretics), decreased cardiac output (heart failure, cardiogenic shock), systemic vasodilation (sepsis, anaphylaxis, liver failure)
  \item \textbf{Intrinsic renal:} Acute tubular necrosis (ischemic or nephrotoxic), acute glomerulonephritis, acute interstitial nephritis (drug-induced), rhabdomyolysis, hemolytic uremic syndrome
  \item \textbf{Post-renal:} Benign prostatic hyperplasia, prostate cancer, ureteral stones, retroperitoneal fibrosis, neurogenic bladder, bilateral ureteral obstruction
\end{itemize}

\section*{Indications for Evaluation}
Seek care if:
\begin{itemize}
  \item Complete anuria (no urine output for 12+ hours), suggesting bilateral obstruction or cortical necrosis
  \item Signs of uremia: nausea, confusion, pericarditis, asterixis, or bleeding
  \item Hyperkalemia on ECG (peaked T waves, widened QRS) or symptomatic fluid overload (pulmonary edema)
\end{itemize}

\section*{Related Symptoms}
\begin{itemize}
  \item Edema, weight gain, hypertension, jugular venous distension
  \item Confusion, fatigue, pruritus, metallic taste (uremic symptoms)
  \item Flank pain or suprapubic discomfort
\end{itemize}
\textbf{Note:} Distinguishing prerenal from intrinsic renal azotemia using the fractional excretion of sodium (FENa) and urine microscopy is essential.

\section*{Self-Care / Home Management}
\begin{itemize}
  \item Avoid NSAIDs, ACE inhibitors, and other nephrotoxic agents during volume depletion
  \item Maintain adequate oral hydration in mild volume depletion; do not force fluids if oliguria persists
  \item Monitor daily weights and urine output
  \item Do not use potassium-containing salt substitutes or electrolyte sports drinks without guidance
\end{itemize}

\section*{Diagnostic Approach}
\begin{itemize}
  \item Bladder catheterization to confirm true oliguria and relieve outlet obstruction
  \item Serum creatinine, BUN, electrolytes, and urinalysis with microscopy (muddy brown casts suggest ATN; dysmorphic RBCs suggest glomerulonephritis)
  \item Renal ultrasound to assess for hydronephrosis, kidney size, and cortical thickness
  \item Urine electrolytes (FENa, FEurea) and renal duplex if vascular cause is suspected
\end{itemize}

\section*{Treatment / Management Overview}
\begin{itemize}
  \item Prerenal: volume resuscitation with isotonic crystalloid, treat underlying cause (pressors for sepsis, inotropes for cardiogenic shock)
  \item Intrinsic: discontinue nephrotoxins, optimize hemodynamics; loop diuretics may manage volume overload but do not alter recovery
  \item Post-renal: bladder catheterization or ureteral stenting/nephrostomy to relieve obstruction
  \item Renal replacement therapy (dialysis) for refractory hyperkalemia, acidosis, volume overload, or uremia
\end{itemize}

\clearpage
